\subsection{The complete count for Liouville integrability}
\label{sec:canonical-independence}

The spectral invariants are complete when the phase coefficients are
distinct within each group. Equal phases allow additional continuous
changes of mode basis, whose conserved quantities supply any missing
integrals. We state the complete count here; Appendix~\ref{sec:integrability-proof}
gives its independence proof for arbitrary phase coefficients.

Write $s_h=D_{hh}$ and $e_h=X_{hh}$, $1\leq h\leq d$.
Partition each original group into maximal sets $B$ with a common
intercept $e_h$. We call these sets equal-phase blocks; their slopes
$s_h$ are also constant, with common value $s_B$.
Define the block moment matrix
\begin{equation}\label{eq:canonical-block-moment}
 (\mathsf J_B)_{ij}
 =\sum_{\substack{h=1\\s_h\ne s_B}}^d
       \frac{X_{ih}X_{hj}}{s_B-s_h},
 \qquad i,j\in B.
\end{equation}
For a constant $|B|\times|B|$ matrix $\mathsf Z$, let $\widehat{\mathsf Z}$ denote its
embedding by zeros outside $B$. Then
\[
 \{\mathsf X,\tr(\mathsf Z\mathsf J_B)\}=-i[\widehat{\mathsf Z},\mathsf X].
\]
Thus its Hamiltonian flow is conjugation by $e^{-i\epsilon\widehat{\mathsf Z}}$.
It preserves the Hamiltonian, since $\widehat{\mathsf Z}$ commutes with the fixed phase
matrices, so the entries of $\mathsf J_B$ are Noether integrals.
Different blocks act independently. Within each block the moment
entries have $-i$ times the matrix Lie--Poisson bracket.
These are the standard quadratic moment maps for rectangular canonical
matrix pairs \cite[Section~2]{Harnad1994}.

Fix an ordering of the indices in $B$, and let $\mathsf J_B^{[t]}$
denote its leading $t\times t$ principal submatrix. The functions
\begin{equation}\label{eq:canonical-moment-chain}
 I_{B,t,k}=\frac1k\tr\!\left((\mathsf J_B^{[t]})^k\right),
 \qquad 1\leq k\leq t\leq |B|,
\end{equation}
commute. For two nested principal blocks, the gradient of an
invariant of the larger block commutes with that block; the
Lie--Poisson pairing with the smaller-block gradient is therefore zero.
This is the Gelfand--Zeitlin construction
\cite{GuilleminSternberg1983}. These functions also commute with every
spectral coefficient, since block conjugation preserves $F$.
Let $\mathcal A$ be the polynomial algebra generated by these functions
and the coefficients of $F$. It is a commuting algebra of first
integrals of $H$.

\begin{theorem}\label{thm:canonical-independence}
For every positive $L,M,N$ and every fixed choice of phase coefficients,
$\mathcal A$ has generic differential rank $p=LM+MN+NL$.
Thus \cref{eq:canonical-H} is generically Liouville integrable.
A complete local set can include $H$ wherever $\mathrm dH\ne0$.
\end{theorem}
The proof in Appendix~\ref{sec:integrability-proof} proceeds by
increasing one matrix dimension at a time. At each insertion, the
spectral coefficients and principal moments together supply one new
independent integral per new canonical pair. This includes arbitrary
coincident phase coefficients. The same argument quantifies the
spectral deficit in \cref{eq:canonical-spectral-deficit}.

When all phases within each original group are distinct, the spectral
family alone is complete; an explicit selection by
line intersections and a spanning tree is given in
\cref{sec:distinct-phase-integrals}. Retaining the $d-1$ power charges
leaves
\begin{equation}\label{eq:canonical-genus-preview}
 G=p-d+1=LM+MN+NL-L-M-N+1
\end{equation}
spectral integrals on the regular locus used in
\cref{sec:spectral-geometry}.

Bolsinov's argument-shift completeness criterion
\cite[Theorem~1.1]{Bolsinov1992} includes the constant bracket
\cref{eq:canonical-frozen}. Its conditions are that
$\mathsf X+\nu\mathsf D$ is regular for every finite $\nu$ and that
the projected kernel has the required dimension; the latter becomes
distinctness of the phases within each original group.
A smooth spectral curve supplies the first condition, so completeness
in the distinct-phase case is a known consequence of this criterion.
\Cref{thm:canonical-independence} gives a uniform direct proof and
completes the count for repeated phases using the established
moment-map and Gelfand--Zeitlin constructions.
Jiang and Rauch-Wojciechowski prove completeness for row counts two
and three (our $L+M$), and leave four rows unresolved
\cite[Sections II.B--II.C]{JRW1991}. The insertion proof covers
arbitrary row and column counts and every fixed phase profile.
